\scanpage{45}
\uline{Relations} a) Morphismes $X\longrightarrow\unsure{\check{\check{P}}|\check{P}}$ et syst. lin. de div \unsure{sans pts fixes paramétrés} par $P$.

\deleted{Applications}
b) Syst. lin. de div. ``\unsure{str.}'' sans ``p. fixe'' paramétré par $\check{P}$, et applications rationnelles sur $S$
\[
X\longrightarrow\check{P}
\]
avec \unsure{petite} propriété supplémentaire \uline{de nature locale}
\margin{[qui est inutile si $X$ est régulier \deleted{\illegible{1}} ou strict loc. factoriel]}

c)
\[
\underline{\mathrm{Syst.lin}}_{X/S}\xrightarrow{\lambda}\underline{\mathrm{Pic}}_{X/S}
\]
morphisme naturel : (\unsure{ayant} \unsure{d'ailleurs} une section canonique). Si $X/S$ est propre à fibres géom. intègres, alors $\lambda$ est rel. représentable comme la grassmannienne complète d'un schéma de Brauer-Séveri généralisé (Attention, il y a une question d'effectivité des données de descente dans un fibré projectif relatif un Module non loc. libre !].

Si, au lieu de fibres géom. intègres, on suppose $f_*(O_X)\simeq O_S$ \unsure{univ\textsuperscript{t}}, on trouve encore un morphisme représentable par un \uline{ouvert} d'un fibré \deleted{projectif} de Brauer-Séveri généralisé .....
\margin{\uline{\unsure{à vérifier}}\par
\unsure{Montrons|Montrer}\par
$X\to S$ est \uline{projectif}\par
si $S$ hensélien local à fibres \unsure{géoréd!}\par
si $\unsure{q}(\unsure{\mathrm{Pic}}(X/S))$\par
est \uline{loc} \unsure{noéthérien}}
\mspage{44}
